\documentclass[a4paper,12pt,fleqn]{scrartcl}
\usepackage[T1]{fontenc}
\usepackage{ucs}
\usepackage[utf8x]{inputenc}
\usepackage{ngerman}
\usepackage[ngerman]{babel}
\usepackage{lastpage}
\usepackage[pdftex]{color,graphicx}

\usepackage[inner=2cm,outer=2cm,top=1cm,bottom=1.5cm,includeheadfoot]{geometry}
\usepackage{fancyhdr}

%---- Title Page
\title{Wochenbericht KW 17}
\author{Christoph Gnip}
\date{ 22.4.-28.4.2013 }

%---- End Title Page

%---- PageLayout
\pagestyle{fancy}

\setlength{\headsep}{10mm}

\usepackage{eso-pic}

% HEADER -------------------------------------------------------------
\fancyhead[R]{
  \includegraphics[width=100pt,keepaspectratio]{img/amedo2012.png}
}

\fancyhead[L]{
  \begin{tabular}[b]{l}
  Christoph Gnip\\
  Projekt: RFID-Evolution
  \end{tabular}
}

%Linie oben
\renewcommand{\headrulewidth}{0.5pt}

% HEADER -------------------------------------------------------------
%Fußzeile links bzw. innen
\fancyfoot[L]{Stand: \today}
%Fußzeile rechts bzw. außen
\fancyfoot[R]{\thepage{} von \pageref{LastPage}}
%Linie unten
\renewcommand{\footrulewidth}{0.5pt}

% Import Macros  ----------------------------------------------------------
\newcommand\nn{\newline\newline}

\begin{document}
\setlength{\headheight}{36pt}
%- the Title page --------------------------------------------------------
\begin{titlepage}
% \AddToShipoutPicture*{\BackgroundPic}

% make the title
\maketitle

%this has to be called here, after \maketitle
\thispagestyle{empty}

%----include the time table for this week
%Table generated by Excel2LaTeX from sheet 'Tabelle1'
\label{tbl:TTWeek1}
\begin{table}[right]
    \textbf{Anwesenheit}
    \hfill
    \begin{tabular}{|l|l|l|l|l|l|l|}
    \textbf{Mo} & \textbf{Di} & \textbf{Mi} & \textbf{Do} & \textbf{Fr} &
\textbf{Sa} & \textbf{So} \\

    Bo    & Bo    & Bo    & Ge    & Ge    & x     & x \\
    \end{tabular}%
\end{table}%



\end{titlepage}
%-------------------------------------------------------------------------------
% the document
%-------------------------------------------------------------------------------
%- Section 1 -------------------------------------------------------------------
\section[Allgemeines]{Allgemeines}
Arbeiten an der Masterarbeit aufgenommen.
\subsection{Organisation der Arbeit}
\begin{table}[right]
    \renewcommand{\arraystretch}{1.5}
    \begin{tabular}{lp{11cm}}
      Titel der Arbeit:  [DE] & Entwicklung eines Systems zur Entfernungsabschätzung für Phasen basiertes UHF RFID Tracking durch Verwendung evolutionärer Berechnungsverfahren \\
      Titel der Arbeit:  [EN] & Development of a Distance Estimation System for Phase-Based UHF RFID Tracking by Utilizing Methodes of Evolutionary Computation \\
      Interne Projektbezeichnung: & PRPS-Evo      \\
      Zeitraum: & 22.April-2.September (20 Wochen)
    \end{tabular}
\end{table}
%
Am Donnerstag, den 18.4., fand ein Treffen mit Herrn Prof.
Dr. Frank Bärmann statt. Ihm wurde die Masterthesis vorgestellt und er erklärt
sich als Erstbetreuer für die Arbeit einverstanden.
\newline
%
Die Anmeldung erfolgt am 25.4.
\newline
%
Aufsetzen der in diesem Projekt verwendeten Entwicklungsumgebung(en). Die Entwicklungsumgebung, die im Rahmen dieser Arbeit verwendet werden, wurden im Pflichtenheft zu dieser Arbeit vorgestellt und diskutiert.
\newline
%
Folgende Umgebungen wurden in dieser Woche aufgesetzt:
\begin{itemize}
  \item Installation der virtuellen Maschine, die als Produktivumgebung in diesem
Projekt verwendet wird, als Betriebssystem wird Ubuntu 12.04 verwendet
  \item Übersetzung und Installation der Shark-Library auf der frischen
  Linux-Maschine
  \item Für die Erstellung der Dokumentation und der Thesis selbst, wird \LaTeX{} verwendet.
  Dazu wurde die Umgebung aufgesetzt. Sie steht sowohl auf der Linux-Maschine als auch auf dem Windowshost zur Verfügung.
\end{itemize}

%- Section 2 -------------------------------------------------------------------
\section[Fortschritte]{Fortschritte}
\begin{enumerate}
 \item Einarbeitung in den CMA-ES Algorithmus. Der Algorithmus wurde erfolgreich in Matlab
      implementiert und erste Versuche (Benchmarks) wurden mit verschiedenen Populationsgrößen durchgeführt. Die Ergebnisse des Benchmarks sind in
      Ein Auszug aus dem Quellcode ist im  gezeigt.
 \item
        dasdasdasd
\end{enumerate}

Bereits in dieser Woche ist es gelungen die Beispiel-Programme der Shark-Library zu erstellen.

%- Section 3 -------------------------------------------------------------------
\section[Probleme]{Probleme}
Die Einrichtung der \LaTeX{} Umgebung dauerte länger als erwartet. Es war ein Tag Arbeit
vorgesehen, aber 2 benötigt. \newline
Unter Windows\copyright gab es Stabilitätsprobleme mit dem verwendeten Editor. Besonders Instabilitäten mit der Verwendung von BibTex.

%- Appendix --------------------------------------------------------------------
\begin{appendix}

\newpage
\huge{Anhänge}
\normalsize
\thispagestyle{plain}
\section{Verwendete Abkürzungen}
\dots{}

\newpage
\thispagestyle{plain}
\section{Matlab CMA-ES Code}
\label{lst:cmaes-mat-code}[blablabla]

\newpage
\thispagestyle{plain}
\section{Tabellen}
  \label{tbl:population-compare}

%   \begin{table}[position specifier]
    \begin{table}[h]
      \centering

    \begin{tabular}{lll}
	\hline
	\textbf{Populationsgröße} & \textbf{Ausführungszeit} & \textbf{Funktion}\\\hline
	x & x & x \\
	y & y & y \\

      \end{tabular}
    \caption{Performancevergleich für unterschiedliche Populationsgrößen}
    \end{table}

\end{appendix}


\end{document}